% А4 = 210мм x 297мм.
% 1in = 25.4mm.
\hoffset=4.6mm          % Левая граница 3см = \in + \hoffset.
\textwidth=170mm        % Ширина текста 170см = 210мм - 30мм - 10мм.
\voffset=-5.4mm         % Верхняя граница 2см = \in + \voffset.
\textheight=257mm       % Высота текста 257мм = 297мм - 20мм - 20мм.
\oddsidemargin=0mm
\evensidemargin=0mm
\topmargin=0mm          % Верхний колонтитул.
\headheight=0mm         % Верхний колонтитул.
\headsep=0mm            % Верхний колонтитул.
\footskip=4.5mm        % Отступ нижнего колонтитула.
\footnotesep=0mm
\sloppy
\raggedbottom
\nonfrenchspacing
\righthyphenmin=2
%\pagestyle{plain}      class defined
%\markboth{}{}
\setcounter {page} {2}  %% 3, 7

\parindent=7.5mm        % Абзацный отступ.

%\renewcommand{\baselinestretch}{1}

\usepackage{mdwlist}    % Многострочные метки в листах

% \usepackage{fancyvrb} % TODO: check
% \usepackage{fvextra} % TODO: check
\usepackage{listings} % TODO: check
\lstset{breaklines=true} % TODO: check
% \usepackage{spverbatim} % TODO: check

\usepackage{caption}
\DeclareCaptionLabelSeparator{colon}{~--~}                      % Разделитель в подписях.
\DeclareCaptionJustification{raggedright}{\raggedright}         % Выравнивание по левому краю.
\DeclareCaptionLabelFormat{tableLabelFormat}{\KZtablename\ #2}  % Подпись к таблицам.
\DeclareCaptionLabelFormat{figureLabelFormat}{Figure #2}       % Подпись к рисункам.
\DeclareCaptionStyle{tableStyle}{labelformat=tableLabelFormat, justification=raggedright, font=large, aboveskip=0ex}   % Стиль таблиц.
\DeclareCaptionStyle{figureStyle}{labelformat=figureLabelFormat, justification=centering, font=large, aboveskip=13pt}   % Стиль рисунков.
%\captionsetup[table]{position=bottom}

\large
% управление размещением плавающих объектов
\setcounter {topnumber} {3}
\setcounter {bottomnumber} {3}
\setcounter {totalnumber} {4}
\renewcommand {\topfraction} {.9}
\renewcommand {\bottomfraction} {.9}
\renewcommand {\textfraction} {.1}
\renewcommand {\floatpagefraction} {.8}

%\makeatletter
%\@addtoreset {equation} {section}
%\makeatother

\renewcommand {\thechapter} {\arabic{chapter}}
\renewcommand {\thesection} {\arabic{chapter}.\arabic{section}}
\renewcommand {\thesubsection} {\arabic{chapter}.\arabic{section}.\arabic{subsection}}
\renewcommand {\thesubsubsection} {\alph{subsubsection})}

\setcounter {secnumdepth} {3}
\setcounter {tocdepth} {2}

% для раздвижки строк в gather, ...
\newlength {\x}
 \x=.4ex plus .1ex minus .1ex

% знак транспонирования
%  \newsavebox {\XXX}
%  \sbox {\XXX} {\hbox{\Large$\scriptscriptstyle T$}}
%\newcommand {\tran} {{\usebox{\XXX}}}

%\def\endproc{\rm} % заканчивает теоремы и т.д.
%\newcommand{\proof}{\par{\bf Доказательство.} }
%\newcommand {\V} {\par\medskip\par} % перед и после теорем, лемм и т.д.

%\newcommand {\rep} [1] {#1\discretionary{}{\hbox{$#1$}}{}}
%\newcommand {\REP} [1] {#1\discretionary{}{\hbox{$#1$}}{}}

%%%%%%%%%%%%%%%%%%%%%%%%%%%%%%%%%%%%%%%%%%%%%%%%%%%%%%%%%%%%%%%%%%%%%

%   Определяем знак полного дифференциала.
\newcommand {\intd}{\mspace{3mu}\mathrm{d}}

%   Переопределяем знаки неравенства согласно сов.стандарту.
\renewcommand {\ge}  {\geqslant}
\renewcommand {\geq} {\geqslant}
\renewcommand {\le}  {\leqslant}
\renewcommand {\leq} {\leqslant}
